% !TeX root = ../main.tex
\chapter{结论与展望}

\section{主要结论}

本文针对复杂地质构造背景下地球物理反演精度不足与泛化能力受限的问题，开展了多尺度特征融合网络构建、物理算子嵌入约束以及不确定性量化三个方面的研究，得出如下主要结论：

\begin{enumerate}[label=（\arabic*）]
  \item 引入跨尺度空洞卷积金字塔与残差注意力机制后，反演网络对薄互层与断裂带等细微地质异常的刻画能力明显增强，有效克服了深层特征提取中的高频衰减现象。
  \item 将可微正演物理算子融入端到端训练闭环，在保持毫秒级推理效率的同时显著降低了反演解的物理失配度，增强了模型在低信噪比与跨工区场景下的泛化稳定性。
  \item 结合蒙特卡洛随机失活与模型集成构建的后验方差估计能够准确指示低置信度反演区域，为地球物理定量解释提供了直观的可靠性度量。
\end{enumerate}

\section{研究局限与未来展望}

尽管本文方法在多个实验场景中取得了理想效果，但仍存在若干值得深入探索的方向：其一，当前可微物理层主要面向二维声波近似，未来可进一步扩展至三维弹性波及各向异性介质情形；其二，如何在多模态地球物理资料（如地震--重磁--电磁联合观测）下构建统一的基础大模型，将是下一步研究的重点内容。
